% chapters/14/rangkuman.tex
\bagian{Rangkuman}

\begin{rangkumanbox}
\begin{itemize}
  \item Aset diolah ketika diimpor, bukan ketika program dijalankan. Hasil
        pengolahannya disimpan, sehingga mengubah \istilah{importer setting}
        mengubah seluruh benda yang memakai aset itu sekaligus.
  \item Empat pengaturan impor yang paling berpengaruh pada aset dua dimensi:
        \istilah{Sprite Mode}, \istilah{Pixels Per Unit}, \istilah{Filter Mode},
        dan \istilah{Compression}.
  \item \istilah{Pixels Per Unit} mengaitkan ukuran gambar dengan ukuran level.
        Nilainya sebaiknya sama untuk seluruh aset, agar ukurannya dapat
        dibandingkan.
  \item Ukuran berkas gambar tidak boleh diubah di luar Unity, sebab
        perubahannya akan menjalar ke seluruh benda yang memakainya.
  \item Komponen \istilah{Animator} adalah mesin keadaan yang mengurus animasi.
        Keadaannya berupa \istilah{animation clip}, dan transisinya diatur
        melalui parameter.
  \item Tiga jenis parameter: angka pecahan, nilai benar atau salah, dan
        pemicu. Pemicu dipakai untuk kejadian yang terjadi sekali, sebab
        mesin mengembalikan nilainya sendiri.
  \item Nama parameter harus sama persis dengan yang ada pada jendela
        \istilah{Animator}; kesamaannya tidak diperiksa oleh penyunting, dan
        kesalahan penamaannya tidak menimbulkan pesan apa pun.
  \item Antarmuka pengguna ditempatkan pada \istilah{Canvas}, sehingga
        susunannya tetap sama pada setiap ukuran layar. Antarmuka hanya
        menampilkan keadaan, sedangkan keadaannya tetap disimpan pada kode.
  \item Penggantian aset dilakukan menurut urutan: pastikan mekaniknya selesai,
        impor seluruh aset, ganti benda satu per satu, lalu periksa bentuk
        batasnya. Bentuk batas tidak berubah otomatis bersama gambarnya.
  \item Aset buatan orang lain hanya boleh dipakai menurut lisensinya, dan
        sumbernya perlu dicatat sejak diimpor.
  \item Bab ini melengkapi bagian visual dari \CPMK{4}. Yang tersisa adalah
        penyajian bunyi pada Bab~XV, dan pengujian menyeluruh pada Bab~XVI.
\end{itemize}
\end{rangkumanbox}
